% /Users/lukaswolff/Desktop/AnsysValidation/TestLukas/informe.tex
\documentclass[11pt,a4paper]{article}
\usepackage[utf8]{inputenc}
\usepackage[T1]{fontenc}
\usepackage[spanish]{babel}
\usepackage{lmodern}
\usepackage{microtype}
\usepackage{geometry}
\geometry{margin=2.5cm}
\usepackage{hyperref}
\hypersetup{colorlinks=true,linkcolor=blue,citecolor=blue,urlcolor=blue}
\usepackage{listings}
\usepackage{xcolor}

\lstset{
    language=Python,
    basicstyle=\ttfamily\small,
    keywordstyle=\color{blue}\bfseries,
    stringstyle=\color{blue},
    commentstyle=\color{gray}\itshape,
    showstringspaces=false,
    breaklines=true,
    frame=single,
    numbers=left,
    numberstyle=\tiny,
    captionpos=b
}

\title{Pyfluent apuntes Lukas W}
\author{ }
\date{\today}

\begin{document}

\maketitle

En primer lugar es nesesario encontrar donde se encuentra la carpeta de fluent, por ejemplo:

\begin{lstlisting}[language=Python]
    "C:\Program Files\ANSYS Inc\ANSYS Student\v252\fluent\ntbin\win64\fluent.exe"
\end{lstlisting}

Es decir, en este caso tengo la version 2025 R2 de ansys student.

Hay que crear una carpeta temporal donde se alojara la sesion de fluent, abrir la terminal de windows (cmd) y escribir:

\begin{lstlisting}[language=Python]
    mkdir C:\Temp
\end{lstlisting}

Luego hay que ejectuar en la terminal lo siguiente:

\begin{lstlisting}[language=Python]
    & "C:\Program Files\ANSYS Inc\ANSYS Student\v252\fluent\ntbin\win64\fluent.exe" 3ddp -g -server -sifile="C:\Temp\fluent-serverinfo.txt"

\end{lstlisting}

Lo cual inicia Fluent 3D sin interfaz grafica y en modo servidor, lo cual deja escuchando para una conexion remota desde python.

Para cerrar Fluent, simplemente hay que escribir en la consola:

\begin{lstlisting}[language=Python]
    exit
\end{lstlisting}

\textbf{Importante:} la version de python a usar debe ser 3.11.9, para eso, se recomienda generar un entorno virtual de python con esa version. Para eso, se recomienda usar Anaconda, preguntar a ChatGPT el paso a paso.

\end{document}